\section{Transfer from the reference volumes}\label{sec:uniform-transfer}

The volume-transfer theorem of \cite[Theorem~A]{VolumeTransfer}
controls products of translated Wick polynomials on the cutoff
tori. We verify its hypotheses for the construction chosen in
\cref{sec:ordered-constants}. The partition comparison, the
nonzero-flux estimate and the conditional-moment argument are
proved in that companion.

The selected constants have $r=2$, $s=1/2$, $M=L^k$ with
$k\ge2$, and $0<\varepsilon\le1/200$. The first inequality
of \eqref{eq:constants-L} implies $L\ge13$.
With $h$ as in \eqref{eq:cutoff-scales}, put
\begin{equation}\label{eq:transfer-scale-choice}
 e_K=e\sqrt h,\qquad v_K=-\log e_K,\qquad
 r_K=\min\{L^i:i\in\Z_{\ge0},\ L^i\ge v_K^2\},
 \qquad R^*=hMr_K/L=hL^{T^*}.
\end{equation}
Write $n_*=R^*/\eta$, $N_*=8L^3$ and
$a_L=(L^3-1)/(L-1)^3$. The positive exponent is
\begin{equation}\label{eq:transfer-charge-exponents}
 \mathfrak e=\tfrac12(1/4-9\varepsilon)-12\varepsilon,
 \qquad \sigma=\mathfrak e-2a_L/N_*>0.
\end{equation}
The terminal inventory in \cref{sec:ordered-constants} gives,
for every $0<e\le e_G$, $e\le1/3$, $v_K\ge8$,
$e_K\log h\le1$, and
\begin{equation}\label{eq:transfer-tail-condition}
 6N_*Mv_K^2\exp(-v_K)
             \bigl(v_K\exp(-v_K)+\tfrac12\bigr)\le\pi^2.
\end{equation}
These are the scale hypotheses of the cited theorem. In particular,
\cite[(41), (49)--(51)]{VolumeTransfer} gives
\begin{equation}\label{eq:transfer-W3}
 6N_*e^2R^*\log(1/e)\le\pi^2,
\end{equation}
and
\begin{equation}\label{eq:transfer-volume-error}
 \bar mR^*>Mr_K/L^2\ge r_K\ge v_K,
 \qquad \exp(-\bar mR^*)\le e\sqrt h.
\end{equation}

Fix the compact real probe $\beta$ and $f=dd^*\beta$.
Let $\delta_{\eta,\mu}$ be the centered difference of step
$\eta$ and $J_\eta$ the cube average. Define
\[
 a'_{\eta,\nu}=-\sum_{\mu\ne\nu}\delta_{\eta,\mu}\beta_{\mu\nu},
 \qquad
 f^\eta_{\mu\nu}=\delta_{\eta,\mu}a'_{\eta,\nu}
                       -\delta_{\eta,\nu}a'_{\eta,\mu}.
\]
The covering source is
$\widehat a^o_\nu(x)=\eta^{1/2}(J_\eta a'_{\eta,\nu})(\eta x+\eta e_\nu/2)$.
It satisfies $d\widehat a^o=(f^\eta)_{[\eta]}$, with the
test normalization in \eqref{eq:limit-test-cover}.
Its reflected, periodized copies $\widehat a_c$ and the
variance $v_\eta=\|d\widehat a^o\|^2$ are those of
\cite[Section~2]{ReferenceResponse}.
Let $D_\beta$ be the side of a least coordinate cube containing
$\supp\beta$, and let $Q_\beta$, centered at $z_\beta$, be
that cube enlarged by two in each direction. Its side is
$D_\beta^{\rm buf}=D_\beta+4$. A centered placement is
$v\in\eta\Z^3$ with
$|v+z_\beta-(\ell_c/2)(1,1,1)|_\infty\le\eta/2$,
where $\ell_c=R^*/(2L)$. These are the conventions of
\cite[Section~2]{ReferenceResponse}. Translate $\widehat a^o$
by $v$ before reflecting and periodizing it to form the cochains
$\widehat a_c$. For such a placement,
Proposition~\extref{11.1}{ReferenceResponse} gives, throughout
$0<e\le e_\sharp(\beta)$,
\[
 D_\beta^{\rm buf}+3\eta\le\ell_c,
 \qquad v+Q_\beta\subset[\ell_c/8,7\ell_c/8]^3,
 \qquad \ell_c=R^*/(2L).
\]
Indeed its placement conditions \exteq{121}{ReferenceResponse}
give both inequalities for $\eta\le1$ and a centered placement.
Thus the full support hypothesis \exteq{6}{VolumeTransfer} holds.

For the same placed probe define the ordinary free derivatives by
\[
 N_0(t)=\frac{\det(D^{\rm ch}[q(t/e)\widehat a_0]+\bar m\eta)}
                   {\det(D^{\rm ch}[0]+\bar m\eta)},
 \qquad c_m=N_0^{(m)}(0),\quad m=2,4.
\]
Write
\[
 H_2(x;v)=x^2-v,\qquad H_4(x;v)=x^4-6vx^2+3v^2,
 \qquad W_m(x;v,c)=e^{-m}H_m(x;v)-c.
\]
The reference response uses the interacting mass
$\boldsymbol\mu_N=\bar m\eta+m_{0,N}(e)$, which can have
either sign; this free ratio uses the positive mass $\bar m\eta$.
Let $\E_{n_*}^{(0)}$ be the all-flux law conditioned on zero
flux, at this same interacting mass, and put
\[
 \Theta_m^{(0)}=
 \left|\E_{n_*}^{(0)}\prod_{c\in\{0,\ldots,2L-1\}^3}
 W_m(\langle d\widehat a_c,u\rangle;v_\eta,c_m)\right|^{1/N_*}.
\]
The exact identity \exteq{53}{ReferenceResponse} identifies
this expression with its normalized differentiated response.
The zero-flux law here and the all-flux cutoff law therefore
have exactly the conventions of \cite{VolumeTransfer}.

Let $z_{r,n}$ denote the twist-averaged zero-flux matter factor
\exteq{1}{TerminalPressure}, with mass $\boldsymbol\mu_N$.
For $N\ge N_{\rm ref}(e)$ and
$0<e\le\min\{e_G,e_\sharp(\beta)\}$, the common-pressure
comparison \exteq{5}{TerminalPressure} and the response
bound \exteq{119}{ReferenceResponse} give fixed constants
$C_{\rm mat},C_{\rm ref}<\infty$ such that
\begin{align}
 L^3\log z_{r,n_*}-\log z_{r,Ln_*}&\le C_{\rm mat},\notag\\
 \Theta_4^{(0)}&\le C_{\rm ref}e^{\mathfrak e},\qquad
 (1+|c_2|)\Theta_2^{(0)}\le C_{\rm ref}e^{\mathfrak e}.
                                      \label{eq:transfer-terminal-input}
\end{align}
Both reference tori use the pair selected on $LR^*$;
their common bulk pressure cancels in this difference.
Theorem~\extref{B}{PositiveProbe} and
Proposition~\extref{5.2}{PositiveProbe}, together with
\eqref{eq:transfer-volume-error}, allow $C_{\rm ref}$ to be
enlarged once so that
\begin{equation}\label{eq:transfer-free-input}
 |c_2|+|c_4|^{1/2}\le C_{\rm ref},\qquad
 |c_4-3c_2^2-b(f)|\le C_{\rm ref}(e+\eta^{1/4}).
\end{equation}
Here the support fit gives $R^*\ge D_\beta+1$.
All constants are independent of charge, mesh and placement.
Choose $C_U$ to be the explicit constant
\exteq{59}{VolumeTransfer}, with these reference constants.

\begin{proposition}\label{prop:transfer-UV-inputs}
For the construction and probe above, throughout
$0<e\le\min\{e_G,e_\sharp(\beta)\}$ and
$N\ge N_{\rm ref}(e)$, there are $u_2,u_4\ge0$, independent
of the outer cutoff $T$, such that
\begin{equation}\label{eq:transfer-small-Wick}
 u_4\le C_Ue^\sigma,\qquad
 (1+|c_2|)u_2\le C_Ue^\sigma,\qquad
 |c_4-3c_2^2-b(f)|\le C_U(e^\sigma+\eta^{1/4}).
\end{equation}
For every $T\ge T^*+1$ and finite
$Y\subset(R^*/L)\Z^3$ of points incongruent modulo $R\Z^3$,
\begin{equation}\label{eq:transfer-translated-Wick}
 \left|\E_{N,T}\prod_{y\in Y}
 W_m(F^\sharp(\tau_yf^\eta);v_\eta,c_m)\right|
                      \le u_m^{|Y|},\qquad m=2,4.
\end{equation}
The fixed smooth test satisfies
\begin{equation}\label{eq:transfer-smooth-cumulant}
 |\kappa_4(F^\sharp(f))-e^4b(f)|
 \le(8C_U+3C_U^2)e^{4+\sigma}
                   +C_U(\eta^2+e^4\eta^{1/4}).
\end{equation}
The tests $f^\eta$ have a common compact support,
$f^\eta\to f$ in $L^2$, and $v_\eta\to\|f\|_2^2$.
\end{proposition}

\begin{proof}
The hypotheses of Theorem~\extref{A}{VolumeTransfer}
have just been checked. Its equations (8)--(10), or
Proposition~\extref{6.1}{VolumeTransfer}, give the assertions.
The carrying-torus condition is satisfied because
$n=8L^{T-T^*}n_*$ is a multiple of $p_{\rm cell}=n_*/L$
and is at least $Ln_*$ when $T\ge T^*+1$.
Although $n_*$ is odd, the reference array has $2L$ cells
per side; its half-integer cell widths are precisely those
used in the cited theorem. Its source identity, test-error
bound and variance convergence are \cite[Section~2.2, (19)--(20)]{VolumeTransfer}.
\end{proof}

The translated-product estimate is the input needed for
components: a bound on the parent's fourth cumulant alone
does not bound the fourth cumulants of its components.
Lemma~\extref{8.1}{VolumeTransfer} bounds the invariant
conditional expectations from these products.
Proposition~\extref{8.2}{VolumeTransfer} then gives the
component cumulant estimate. Theorem~\extref{A}{VolumeTransfer}
performs both steps for any prescribed volume-first parent
whose components satisfy the stated Gaussian bound.
We apply this conclusion to the parent and component set
in the proof of \cref{thm:main}.
